\section{Introduction}\label{sec:introduction}

Topological data analysis (TDA) and causal inference (CI) are the two
mathematical languages in which the modern data-analytic enterprise
speaks most precisely about \emph{shape} and about \emph{consequence}.
TDA, developed since the late 1990s, supplies a multi-scale,
functorial, and provably stable grammar for the geometric structure of
data \citep{edelsbrunner2002topological,carlsson2009topology,
edelsbrunner2010computational,wasserman2018topological,
chazal2021introduction}. Causal inference, developed in parallel since
the 1970s, supplies a complete calculus for the interventional and
counterfactual content of probabilistic models
\citep{rubin1974estimating,pearl1995causal,pearl2009causality,
imbens2015causal,peters2017elements,bareinboim2020pearls}. Both fields
have reached methodological maturity; both have generated their own
journals, their own statistical scaffolding, and their own
implementation ecosystems; and both have been applied to many of the
same empirical objects---networks, dynamical systems, biomedical
images, spatiotemporal fields---without explicitly meeting.

The thesis of this paper is that the absence of an explicit meeting
between the two is no longer tenable. Modern scientific data are
increasingly complex, high-dimensional, and non-Euclidean, in ways
that undermine the parametric and structural assumptions on which
classical causal estimators depend; and topological data analysis,
having absorbed a decade of statistical inference results, is now
positioned to deliver not merely descriptive summaries but
\emph{causal} estimators with calibrated uncertainty. A small but
growing body of work---the four seed papers
\citep{ibeling2021topological,kim2026topological,lecci2014statistical,
elyaagoubi2023statistical} that we will treat as the anchors of the
field, together with the sheaf-theoretic causal-inference programme
\citep{mahadevan2021causal,mahadevan2022universal,mahadevan2025sheaf}
and the topological-ordering family of causal-discovery algorithms
\citep{rolland2022score,tran2024constraining,xu2025information,
nocurl2024,wu2026optimization}---has begun to formalise the
connection. What is missing is the framework that organises these
contributions into a single research domain.

We propose such a framework. The remainder of this introduction
positions the contribution: \Cref{sec:intro_parallel} traces the
parallel evolution of TDA and CI; \Cref{sec:intro_why} argues for
their integration; \Cref{sec:intro_limits} states the limitations of
each tradition in isolation; and \Cref{sec:intro_contributions}
enumerates the four contributions that we will make in the body of
the paper.

\subsection{The parallel evolution of two fields}\label{sec:intro_parallel}

Topological data analysis crystallised as a discipline with the
introduction of \emph{persistent homology}: the assignment, to a
filtered simplicial complex, of a multi-set of birth--death pairs that
record the appearance and disappearance of topological features at
each scale \citep{edelsbrunner2002topological,zomorodian2005computing,
edelsbrunner2008persistent}. The stability theorem of
\citet{cohen2007stability}, refined by \citet{chazal2016structure},
established that persistent homology is $1$-Lipschitz with respect to
the bottleneck distance under $L^\infty$ perturbations of the filter,
turning a previously qualitative summary into a calibrated descriptor.
The decade that followed delivered the statistical inference
programme---confidence sets, bootstrap, rates of convergence, and
functional summaries (landscapes, silhouettes, images)---synthesised
in \citet{lecci2014statistical} and consolidated by
\citet{bubenik2015statistical,adams2017persistence,
chazal2014stochastic}. \citet{carlsson2009topology}, the
\citet{singh2007mapper} Mapper algorithm, and the categorical
formalism of \citet{spivak2014category,curry2014sheaves} situated
TDA within a broader functorial-data-analysis programme that
emphasised composability and structure.

Causal inference, in parallel, developed in two complementary
traditions. \citet{rubin1974estimating} and the potential-outcomes
school \citep{rubin2005causal,imbens2015causal} grounded the field in
counterfactual contrasts and identification conditions
(consistency, ignorability, positivity), and produced the
doubly-robust efficient estimation machinery
\citep{robins1989probability,kennedy2024semiparametric,
chernozhukov2018double} that today underwrites empirical work in
biomedicine, econometrics, and policy science. \citet{pearl1995causal,
pearl2009causality,pearl2018book} introduced the structural-causal-
model formalism and do-calculus, which axiomatised intervention as a
syntactic operation on directed acyclic graphs and enabled
\emph{causal discovery}: the recovery of $G$ itself from data via
constraint-based \citep{spirtes2001causation}, score-based
\citep{meek1995strong}, and functional \citep{shimizu2006lingam,
peters2014causal} methods. \citet{bareinboim2020pearls,
ibeling2021topological} crystallised the causal hierarchy
(observational, interventional, counterfactual) into a single
expressive ladder.

\Cref{fig:timeline} in \Cref{sec:literature} traces the two
trajectories side by side. Until $\sim$2020 there is essentially no
crossing: a researcher could attend the major venues of either field
for a decade without hearing the other's vocabulary. The crossings
that began to appear thereafter---a topology on the space of SCMs
\citep{ibeling2021topological}; topology-aware treatment-effect
estimators \citep{kim2026topological}; simulation-based topological
inference for multivariate time series
\citep{elyaagoubi2023statistical}; sheaf-theoretic causal calculi
\citep{mahadevan2021causal,mahadevan2022universal,mahadevan2025sheaf};
and topological-ordering algorithms for causal discovery
\citep{rolland2022score,tran2024constraining,xu2025information,
nocurl2024,wu2026optimization}---are the methodological seeds of the
field this paper proposes to name.

\subsection{Why topology matters for causality}\label{sec:intro_why}

Three properties of topology recommend it as a natural partner for
causal inference, and each is reflected in one of the four domains of
our framework (\Cref{fig:taxonomy}).

\paragraph{Non-Euclidean data.}
The classical average treatment effect, $\ATE = \E[Y^1 - Y^0]$,
implicitly assumes that the difference $Y^1 - Y^0$ is well defined and
scalar-valued. Modern outcomes routinely fail both assumptions: a CT
scan, a molecular conformation, a brain connectivity graph, and a
gridded climate field are all non-Euclidean objects on which
``$Y^1 - Y^0$'' has no obvious meaning. Topological summaries of such
objects---persistence diagrams, Reeb graphs, Mapper complexes---live
in calibrated metric spaces (the Wasserstein space of diagrams, the
interleaving space of Reeb graphs) on which causal contrasts are
well defined \citep{kim2026topological,kurisu2024geodesic}. Sub-domain
A3 of \Cref{sec:framework} (causal embeddings in persistence diagrams)
makes this construction precise.

\paragraph{Robustness through stability.}
The stability theorem (\Cref{eq:stability}) is the operational core of
the field: small perturbations in the data produce only small
perturbations in the persistent invariants, with a constant of one.
This is precisely the kind of bound that classical causal estimators
typically lack and that practical deployment increasingly demands.
\Cref{prop:stability-transfer} shows that any Lipschitz causal
functional of a persistent invariant inherits the quantitative
robustness of $\PHk{k}$, giving a TCDA estimator a tractable
sensitivity guarantee with respect to Hausdorff-scale perturbations
of the data support. Sub-domain B1 of \Cref{sec:framework} formalises
this transfer.

\paragraph{Hidden structure and confounding.}
Unobserved confounders move observed variables along unobserved
orbits and so deposit a topological signature in the joint support
that no graphical or moment-based diagnostic detects. The
worked example of \Cref{sec:cs_confounding}---a hidden circular
confounder generating a persistent $1$-cycle in the residual support
of an apparently null treatment effect---illustrates the principle,
and \Cref{prop:confounding_signature} states it formally. Sub-domain
B2 of \Cref{sec:framework} positions topological diagnostics as a
complement to the meagre-set theorem of
\citet{ibeling2021topological}: where the theorem tells us that
assumption-free identification is impossible on almost all SCMs,
topological diagnostics tell us \emph{which} SCMs have residual
geometric structure consistent with hidden confounding.

These three properties are sufficient to motivate the framework, but
the case for TCDA also rests on a fourth, less elementary observation:
much of modern causal discovery is \emph{already} doing topological
data analysis, often without naming it as such.
\citet{rolland2022score,tran2024constraining,xu2025information,
nocurl2024,wu2026optimization} all extract a topological order from
the data and then refine edges within that order; the order itself is
an invariant of the joint distribution under the appropriate functor,
and \Cref{prop:identifiability} formalises the corresponding
identifiability question. Naming what is already implicit is a small
contribution, but it is one that the field requires before it can
begin to make explicit progress on the deeper questions.

\subsection{Limitations of current approaches}\label{sec:intro_limits}

The case against the status quo is symmetric. Classical causal
inference, applied to non-Euclidean outcomes, must reduce them to
scalar surrogates and so discards exactly the structure of interest;
classical TDA, applied to data with a known causal-graph context,
exposes geometric features but is silent on the question of which
geometric feature carries an intervention's effect. Neither field, on
its own, is equipped to handle a treatment whose outcome is a tumour
shape, a folded protein, a brain network, or a precipitation cluster.

A further limit is structural rather than technical.
\citet{ibeling2021topological}'s topological causal hierarchy theorem
shows that the set of SCMs on which assumption-free identification of
$P(y \mid \doop(x))$ is possible is \emph{topologically meagre}---in
the Baire-category sense, exceptional. This is the topological
counterpart of measure-theoretic ``no-free-lunch'' results
\citep{oxtoby1971measure,belot2020absolutely,genin2017topology,
genin2018topology}, and it places a precise upper limit on what any
purely observational causal procedure can hope to achieve. The
theorem also identifies the structure that must be supplied to do
better, namely a topology on the model space, and is one of the
clearest demonstrations that the language of CI requires topological
extension on its own terms. The complementary critique runs from CI
back to TDA: a persistence diagram is the same whether the variables
it summarises stand in a causal relationship or are mere coincident
features, so without an intervention semantics overlay, a TDA
descriptor cannot license a causal claim. \Cref{sec:lit_gaps}
formalises five specific gaps that follow from this asymmetric
maturity, and \Cref{sec:research} translates each into a concrete
open problem.

\subsection{Our contributions}\label{sec:intro_contributions}

The paper makes four contributions.

\begin{enumerate}[label=(\arabic*),leftmargin=2em]
  \item \textbf{A formal framework.} We propose \emph{Topological
        Causal Data Analysis} (TCDA) as a coherent research domain at
        the intersection of TDA and CI, give a formal definition of a
        TCDA problem as a tuple $(\Mc, G, \Tc, \Cc)$
        (\Cref{def:tcda}), state two foundational propositions
        (\Cref{prop:stability-transfer,prop:identifiability}) that
        establish well-posedness, and defend a taxonomy of four
        top-level domains and twelve sub-domains
        (\Cref{fig:taxonomy}). This is the central intellectual
        contribution and is the subject of \Cref{sec:framework}.

  \item \textbf{A structured literature review.}
        \Cref{sec:literature} synthesises the existing connections
        between TDA and CI along the taxonomy of
        \Cref{sec:framework}, organised across three eras
        (\Cref{sec:lit_early,sec:lit_method,sec:lit_app}) and
        closing with five concrete gaps (\Cref{sec:lit_gaps}). The
        review establishes that every node of the taxonomy is
        already anchored in published work, and that the missing
        piece is the unifying framework we provide.

  \item \textbf{A suite of proof-of-concept case studies.}
        \Cref{sec:case_studies} instantiates the TCDA tuple in five
        concrete settings---biomedical imaging and molecular
        structure (\Cref{sec:cs_tate}), multivariate brain time
        series (\Cref{sec:cs_cycles}), confounding diagnostics
        (\Cref{sec:cs_confounding}), Mapper-based causal abstraction
        with calibrated uncertainty (\Cref{sec:cs_mapper}), and
        spatiotemporal climate attribution
        (\Cref{sec:cs_climate})---together with a worked
        confounding-signature proposition
        (\Cref{prop:confounding_signature}) and a synthesis table
        (\Cref{tab:cs_summary}) mapping each case study back to the
        taxonomy.

  \item \textbf{A research agenda.} \Cref{sec:research} formalises
        eleven Open Questions (\Cref{q:causal-modules}--%
        \Cref{q:topo-docalculus}), fourteen Research Avenues
        (\Cref{r:causal-mapper}--\Cref{r:benchmarks}), and four
        Computational Challenges (\Cref{c:scalable-ph}--%
        \Cref{c:gpu-cubical}), organised into theoretical,
        methodological, computational, domain-specific, and
        industry-facing branches and summarised in the mind-map of
        \Cref{fig:research_mindmap}. Each item is calibrated to be
        specific enough to underwrite an individual research project
        or grant proposal.
\end{enumerate}

\paragraph{What this paper is not.}
TCDA is a proposal, not a theorem. We do not claim that every
component of the framework is mature, that every open question is
solvable, or that the taxonomy is final. Several of the propositions
in \Cref{sec:framework,sec:case_studies} are stated as proof sketches
rather than full theorems, and several of the open questions in
\Cref{sec:research} are likely to be reformulated as the field
matures. The contribution of the paper is to articulate the
framework, locate the existing literature within it, and identify
the questions whose answers will determine whether TCDA becomes a
durable research domain or remains a useful but provisional
re-organisation of work that already exists.

\paragraph{Organisation.}
\Cref{sec:preliminaries} collects the minimal TDA and CI background
needed to read the rest of the paper.
\Cref{sec:framework}~develops the TCDA framework and proves its
foundational propositions. \Cref{sec:literature} surveys the existing
state of the field. \Cref{sec:case_studies} presents five
proof-of-concept case studies. \Cref{sec:research} states the
research agenda. \Cref{sec:conclusion} closes with a call for
interdisciplinary collaboration.
